

\section{Background}
% In this section, we explain the hardware primitives that \thename incorporates into its design and the fundamentals of capability-based access control.
% Also, we introduce the notations that we will use consistently throughout the paper.

\subsection{Pointer Authentication (PA)}
\label{sec:background:pac}
\emph{\Glsentryfull{pa}}, introduced on ARMv8.3-A~\cite{arm-arch}, provides hardware acceleration and ISA extension for cryptographically authenticated pointers.
The design of \gls{pa} is inspired by previous works that protect pointers from corruption~\cite{ccfi,cpi}.
\gls{pa} allows attaching a \gls{pac} to a pointer in the unused bits [54:48] of a 64-bit address.
A \gls{pac} is generated using one of the five keys stored in the newly added registers that can only be accessed or modified with privileged instructions.
The five keys consist of two keys for signing code/instruction pointers (Instruction-\{A,B\}), another two for signing data (Data-\{A,B\}) pointers, and one general-purpose key (G).
We denote these keys as \key{IA},\key{IB}, \key{DA}, \key{DB} and \key{G}.
\Gls{pa} also introduces \texttt{pac} and \texttt{aut} families of instructions to \emph{sign} (i.e., calculate the \gls{pac} and attach it to the value in the operand register) and \emph{authenticate} a 64-bit value with an \emph{optional} modifier using the key indicated by the instruction name.
If authentication on a pointer succeeds, the authentication code in the pointer is cleared, making the pointer usable.
The \gls{pac} value is corrupted if the check fails, which triggers a segmentation fault when the pointer is dereferenced.
We denote the operations of \texttt{pac} and \texttt{aut} instructions that use the key \key{} to sign and authenticate a given data \emph{D} as \pac{}{\emph{D}}{\emph{mod}} and \aut{}{\emph{D}}{\emph{mod}}.


% \emph{\Glsentryfull{pa}}, introduced on ARMv8.3-A~\cite{arm-arch} and ARMv8-M~\cite{armv8-m}, provides hardware acceleration and ISA extension for cryptographically authenticated pointers.
% The design of \gls{pa} is inspired by previous works that protect pointers from corruption (\ie, \emph{pointer integrity})~\cite{ccfi,cpi}.
% \gls{pa} allows attaching a \gls{pac} to a pointer in the unused bits [54:48] of a 64-bit pointer.
% A \gls{pac} is generated using one of the five keys stored in the newly added registers that can only be accessed or modified with privileged instructions.
% The five keys consist of two keys for signing code/instruction pointers (Instruction-\{A,B\}), another two for singing data (Data-\{A,B\}) pointers, and one general-purpose key (G).
% We denote five keys as \key{IA},\key{IB}, \key{DA}, \key{DB} and \key{G}.
% \noh{Moreover, \gls{pa} introduces \texttt{pac} and \texttt{aut} family of instructions to sign and authenticate pointers.
% The \texttt{pac} family of instructions calculate a hash of a given pointer in the operand register

% PA also introduces \texttt{pac} and aut families of instructions to sign (i.e., calculate the PAC and attach it to
% the input) and verify a 64-bit value with an optional modifier using the key indicated by the instruction name. If authentication on a pointer succeeds, the authentication code in the pointer is cleared, making the pointer usable.
% The PAC value is corrupted if the check fails, which triggers a segmentation
% fault when the pointer is dereferenced.

% These instructions optionally take context along with the pointer as operands.

% Intuitively, the signing context associates a given pointer with a context, such that a signed pointer is only valid if its authenticating context is identical.
% The upper bits of a pointer are corrupted if \gls{pa} authentication fail, which will trigger a segmentation fault when dereferenced.
% Hereinafer, we denote the signing and authentication operations that use the key \key{} to sign or authenticate a given pointer \emph{P} as \pac{}{\emph{P}}{\emph{mod}} or \aut{}{\emph{P}}{\emph{mod}}.
% }


\subsection{Memory Tagging Extension (MTE)}
\label{sec:background:mte}
\emph{\Glsentryfull{mte}}~\cite{mte} is another hardware-backed security feature to ARMv8.5-A~\cite{arm-arch}.
\Gls{mte} adapts the principles of tagged memory in the existing research proposals and implementations in other architectures~\cite{timber-v,hw-tagged-memory,sparc-adi}. 
It facilitates memory safety by enabling the 4-bit tagging of pointers and 16-byte aligned address ranges.
When \gls{mte} is enabled, the processor raises an interrupt if the tags differ between the memory and the pointer, allowing vulnerabilities such as buffer overflow to be captured.
We denote \tagz{}{\emph{P}} as the tagging operation that first assigns the 4-bit tag \textbf{T} to the bits [59:56] of a pointer \emph{P}, then tag 16-byte-aligned memory region between \emph{P} and \emph{P}$+$size with \textbf{T}.
In addition, \gls{pa} and \gls{mte} can be \emph{simultaneously} enabled.
Hence, a pointer can be tagged \emph{then} signed, and the resulting pointer would carry a tag in bits [60:56] and \gls{pac} in [54:48].
% Hence, the security guarantees of \gls{pa} and \gls{mte} are complementary.
% NOTE: Trimmed
% \Gls{mte} alone cannot guarantee memory isolation due to the lack of a secure pointer tagging mechanism; an attacker can craft maliciously tagged pointers (\eg, through memory corruption) to access tagged memory of any color.
% On the other hand, \gls{pa} alone is incapable of ensuring memory isolation; an adversary can authenticate a pointer and then perform an out-of-bounds memory (\eg, \texttt{gets(AUT(\emph{P}))}).
% Hence, \thename adapts tagged and signed pointers and turns them into non-forgeable memory references, which we will explain in the following sections.



% \subsection{Capability-based access control}
% \label{sec:background:capability}
% \deleted{
% Capability-based access control revolves around non-forgeable \emph{capability tokens}, or \emph{references}, representing a subject's access rights to protected resources.
% Such capability-based security allows precisely defined \emph{least privilege} and privilege sharing through delegation in compartments.
% Due to its potential in and hence has been sought after by many historical and recent works~\cite{obj-cap,cap-addressing,eros,capsicum,cheri,cheri-risc,wasi,fuchsia,framer}.
% By requiring the capabilities of each software component to be explicitly given, capability-based access control is the ideal scheme for achieving the \emph{least privilege principle} and eliminating \emph{ambient authority} of software~\cite{cheri,capsicum,eros}.
% Access control decisions are made simply by verifying and authenticating the requested subjects' resource references in such systems.
% Notably, recent undertakings towards a design of access control systems from scratch~\cite{fuchsia,wasi}, such as Google's Fushcia OS~\cite{fuchsia}, adapted the capability model.
% }
